\documentclass[../../../include/open-logic-section]{subfiles}

\begin{document}

\olfileid[pt]{sth}{ordinals}{wo}
\olsection{Boas Ordenações}

A noção fundamental é a seguinte:

\begin{defn}
A relação $<$ \emph{bem ordena} $A$ sse cumpre estas duas condições:
\begin{enumerate}
	\item $<$ é conexa, isto é, para todos $a, b \in A$, ou $a < b$
	ou $a = b$ ou $b < a$;
	\item todo o subconjunto não vazio de $A$ tem um !!{element}
	$<$-minimal, isto é, se $\emptyset \neq X \subseteq A$ então $(\exists m \in
	X)(\forall z \in X)z \nless m$
\end{enumerate}
\end{defn}

É fácil ver que os três exemplos que acabamos de considerar eram de fato
relações de boa ordenação.

\begin{prob}
\Olref[sth][ordinals][idea]{sec} apresentou três ordenações exemplo
sobre os números naturais. Verifique que cada uma é uma boa ordenação.
\end{prob}

Eis algumas observações elementares mas extremamente importantes sobre
boa ordenação.

\begin{prop}\ollabel{wo:strictorder}
Se $<$ bem ordena $A$, então todo o subconjunto não vazio de $A$ tem um
único membro $<$-mínimo, e $<$ é irreflexiva, assimétrica e transitiva.
\end{prop}

\begin{proof}
Se $X$ é um subconjunto não vazio de $A$, tem um !!{element} $<$-minimal
$m$, isto é, $(\forall z \in X)z \nless m$. Como $<$ é conexa,
$(\forall z \in X)m \leq z$. Logo $m$ é o !!{element} $<$-mínimo de $X$.

Para a irreflexividade, fixe $a \in A$; o !!{element} $<$-mínimo de $\{a\}$ é $a$, logo $a \nless a$. Para a transitividade, se $a < b < c$, então
como $\{a, b, c\}$ tem um !!{element} $<$-mínimo, $a < c$. A assimetria
segue da irreflexividade e da transitividade
\end{proof}

\begin{prop}\ollabel{propwoinduction}
Se $<$ bem ordena $A$, então para qualquer fórmula $\phi(x)$:
\[
	\text{se }(\forall a \in A)((\forall b < a)\phi(b) \lif 
		\phi(a))\text{, então }(\forall a \in A)\phi(a).
\]
\end{prop}

\begin{proof}
Provaremos a contrapositiva. Suponha $\lnot(\forall a \in
A)\phi(a)$, isto é, que $X = \Setabs{x \in A}{\lnot\phi(x)} \neq
\emptyset$. Então $X$ tem um !!{element} $<$-minimal, $a$. Logo $(\forall
b < a)\phi(b)$ mas $\lnot \phi(a)$. \end{proof}\noindent Esta última
propriedade deve recordar-lhe o princípio de indução forte nos naturais,
isto é: se $(\forall n \in \omega)((\forall m < n)\phi(m)
\lif \phi(n))$, então $(\forall n \in \omega)\phi(n)$. E esta
propriedade torna a boa ordenação numa noção muito \emph{robusta}.\footnote{Lembrete: todas as fórmulas podem ter parâmetros (salvo indicação em contrário).} 

\end{document}
